% GENERATED FILE. Edit canonical lessons, then run npm run generate:books.
% canonical-source-hash: fnv1a64:ba475686c68e1def
% canonical-lessons: BN-C278-dhakka, BN-C278-chola, BN-C278-bhanga, BN-C278-ultano, BN-C278-ahoto

\chapter{A Crash on the Road}
\label{ch:bn-a-crash-on-the-road}
\hlchaptermodality{278}

\begin{chapteropening}
\textbf{By the end of this chapter:} \emph{I can tell what happened in a road accident.}

Complete the last lesson of chapter 278: I can tell what happened in a road accident.
\end{chapteropening}

\section[\texorpdfstring{\bn{ধাক্কা} (dhākkā)}{ধাক্কা (dhākkā)}]{\bn{ধাক্কা} (dhākkā) — a jolt; a collision}

\label{lesson:BN-C278-dhakka}



\begin{quote}
Before the new one: say the Bengali for a crowd, then the Bengali for to turn up.
\end{quote}

\subsection*{You'll want to know: \bn{ধাক্কা}}
\textbf{\bn{ধাক্কা}} — \emph{dhākkā} — \textquotedblleft{}a jolt; a collision\textquotedblright{}.

\begin{quote}
\bn{আস্তে} \bn{না} \bn{চালালে} \bn{ধাক্কা} \bn{লাগবে।}
\end{quote}

\begin{quote}
\emph{āste nā chālāle dhākkā lāgbe.}
\end{quote}

\begin{quote}
If you don't drive slowly, you will crash into something.
\end{quote}

\bn{ধাক্কা} \bn{লাগা}, to collide; \bn{ধাক্কা} \bn{দেওয়া}, to push or bump into.

\begin{cousinweb}
\textbf{\bn{ধাক্কা}} (\emph{dhākkā}) is filed in CDIAL 6701 under a Middle Indo-Aryan \emph{dhakk-}, 'push, strike', which lists Bengali \emph{dhākkā}, 'a push'.
\end{cousinweb}

\begin{grammarlens}[title={What you've built}]
One more word for telling what happened.
\end{grammarlens}

\subsection*{Your turn}
\begin{itemize}
\raggedright
  \item \emph{Say it:} \emph{dhākkā}
  \item \emph{Say it:} \emph{dhākkā}, once more
  \item \emph{Say it:} \emph{āste nā chālāle dhākkā lāgbe.} — If you don't drive slowly, you will crash into something.
  \item \emph{Recall:} say the Bengali for a crowd, then the Bengali for to turn up, then say \emph{dhākkā} again
\end{itemize}

\begin{tcolorbox}[breakable,colback=teal!4,colframe=teal!35!black,title={Before you move on}]
What does \bn{ধাক্কা} mean? (\textquotedblleft{}A jolt; a collision\textquotedblright{}.) Say it once more.
\end{tcolorbox}
\section[\texorpdfstring{\bn{চলা} (chôlā)}{চলা (chôlā)}]{\bn{চলা} (chôlā) — to move; to go along}

\label{lesson:BN-C278-chola}



\begin{quote}
Before the new one: say the Bengali for to turn up, then the Bengali for a jolt.
\end{quote}

\subsection*{You'll want to know: \bn{চলা}}
\textbf{\bn{চলা}} — \emph{chôlā} — \textquotedblleft{}to move; to go along\textquotedblright{}.

\begin{quote}
\bn{বাসটা} \bn{দ্রুত} \bn{চলছিল}, \bn{তাই} \bn{হঠাৎ} \bn{থামতে} \bn{পারল} \bn{না।}
\end{quote}

\begin{quote}
\emph{bāsṭā drutô chôlchhilô, tāi hôṭhāt thāmte pārlô nā.}
\end{quote}

\begin{quote}
The bus was going fast, so it couldn't stop suddenly.
\end{quote}

Past (he/she …): \bn{সেই} \bn{বাস} \bn{চলল}, \emph{sei bās chôllô}.

\begin{cousinweb}
\textbf{\bn{চলা}} (\emph{chôlā}) comes down from Sanskrit \emph{calati}, 'goes away; shakes' (CDIAL 4715, which lists Bengali \emph{calā}, 'to move, walk'). It is the plain partner of \textbf{\bn{চালানো}}: a driver \bn{চালায়}, and the car \bn{চলে}.
\end{cousinweb}

\begin{grammarlens}[title={What you've built}]
One more word for telling what happened.
\end{grammarlens}

\subsection*{Your turn}
\begin{itemize}
\raggedright
  \item \emph{Say it:} \emph{chôlā}
  \item \emph{Say it:} \emph{chôlā}, once more
  \item \emph{Say it:} \emph{bāsṭā drutô chôlchhilô, tāi hôṭhāt thāmte pārlô nā.} — The bus was going fast, so it couldn't stop suddenly.
  \item \emph{Recall:} say the Bengali for to turn up, then the Bengali for a jolt, then say \emph{chôlā} again
\end{itemize}

\begin{tcolorbox}[breakable,colback=teal!4,colframe=teal!35!black,title={Before you move on}]
What does \bn{চলা} mean? (\textquotedblleft{}To move; to go along\textquotedblright{}.) Say it once more.
\end{tcolorbox}
\section[\texorpdfstring{\bn{ভাঙা} (bhāngā)}{ভাঙা (bhāngā)}]{\bn{ভাঙা} (bhāngā) — to break}

\label{lesson:BN-C278-bhanga}



\begin{quote}
Before the new one: say the Bengali for a jolt, then the Bengali for to move.
\end{quote}

\subsection*{You'll want to know: \bn{ভাঙা}}
\textbf{\bn{ভাঙা}} — \emph{bhāngā} — \textquotedblleft{}to break\textquotedblright{}.

\begin{quote}
\bn{ধাক্কায়} \bn{গাড়ির} \bn{জানালা} \bn{ভাঙল।}
\end{quote}

\begin{quote}
\emph{dhākkāy gāṛir jānālā bhānglô.}
\end{quote}

\begin{quote}
The car window broke in the collision.
\end{quote}

Past (he/she …): \bn{সেই} \bn{জানালা} \bn{ভাঙল}, \emph{sei jānālā bhānglô}.

It works both ways: \bn{কেউ} \bn{জানালা} \bn{ভাঙল}, someone broke the window; \bn{জানালা} \bn{ভাঙল}, the window broke. The -\bn{এ} form is \bn{ভেঙে}.

\begin{cousinweb}
\textbf{\bn{ভাঙা}} (\emph{bhāngā}) comes down from Sanskrit \emph{bhaṅga}, 'breaking; a break' (CDIAL 9353, which lists Bengali \emph{bhāṅā}, 'to be broken'). Its \bn{ঙ} is the nasal \emph{ng} before the vowel.
\end{cousinweb}

\begin{grammarlens}[title={What you've built}]
One more word for telling what happened.
\end{grammarlens}

\subsection*{Your turn}
\begin{itemize}
\raggedright
  \item \emph{Say it:} \emph{bhāngā}
  \item \emph{Say it:} \emph{bhāngā}, once more
  \item \emph{Say it:} \emph{dhākkāy gāṛir jānālā bhānglô.} — The car window broke in the collision.
  \item \emph{Recall:} say the Bengali for a jolt, then the Bengali for to move, then say \emph{bhāngā} again
\end{itemize}

\begin{tcolorbox}[breakable,colback=teal!4,colframe=teal!35!black,title={Before you move on}]
What does \bn{ভাঙা} mean? (\textquotedblleft{}To break\textquotedblright{}.) Say it once more.
\end{tcolorbox}
\section[\texorpdfstring{\bn{উল্টানো} (ulṭāno)}{উল্টানো (ulṭāno)}]{\bn{উল্টানো} (ulṭāno) — to overturn; to turn over}

\label{lesson:BN-C278-ultano}



\begin{quote}
Before the new one: say the Bengali for to move, then the Bengali for to break.
\end{quote}

\subsection*{You'll want to know: \bn{উল্টানো}}
\textbf{\bn{উল্টানো}} — \emph{ulṭāno} — \textquotedblleft{}to overturn; to turn over\textquotedblright{}.

\begin{quote}
\bn{বাসের} \bn{ধাক্কায়} \bn{রিকশাটা} \bn{উল্টে} \bn{গেল।}
\end{quote}

\begin{quote}
\emph{bāser dhākkāy rikshāṭā ulṭe gelô.}
\end{quote}

\begin{quote}
The rickshaw overturned when the bus hit it.
\end{quote}

Past (he/she …): \bn{তিনি} \bn{বইয়ের} \bn{পাতা} \bn{উল্টালেন}, \emph{tini bôiyer pātā ulṭālen}.

\begin{cousinweb}
\textbf{\bn{উল্টানো}} (\emph{ulṭāno}) is filed in CDIAL 2368 under a Middle Indo-Aryan \emph{ullaṭati}, 'turns over', which lists Bengali \emph{ulṭāna}, 'to turn over'. \bn{উল্টে} \bn{যাওয়া}, to tip over; \bn{উল্টানো} on its own, to turn something over.
\end{cousinweb}

\begin{grammarlens}[title={What you've built}]
One more word for telling what happened.
\end{grammarlens}

\subsection*{Your turn}
\begin{itemize}
\raggedright
  \item \emph{Say it:} \emph{ulṭāno}
  \item \emph{Say it:} \emph{ulṭāno}, once more
  \item \emph{Say it:} \emph{bāser dhākkāy rikshāṭā ulṭe gelô.} — The rickshaw overturned when the bus hit it.
  \item \emph{Recall:} say the Bengali for to move, then the Bengali for to break, then say \emph{ulṭāno} again
\end{itemize}

\begin{tcolorbox}[breakable,colback=teal!4,colframe=teal!35!black,title={Before you move on}]
What does \bn{উল্টানো} mean? (\textquotedblleft{}To overturn; to turn over\textquotedblright{}.) Say it once more.
\end{tcolorbox}
\section[\texorpdfstring{\bn{আহত} (āhôtô)}{আহত (āhôtô)}]{\bn{আহত} (āhôtô) — injured; hurt}

\label{lesson:BN-C278-ahoto}



\begin{quote}
Before the new one: say the Bengali for to break, then the Bengali for to overturn.
\end{quote}

\subsection*{You'll want to know: \bn{আহত}}
\textbf{\bn{আহত}} — \emph{āhôtô} — \textquotedblleft{}injured; hurt\textquotedblright{}.

\begin{quote}
\bn{দুর্ঘটনায়} \bn{অনেক} \bn{মানুষ} \bn{আহত} \bn{হল।}
\end{quote}

\begin{quote}
\emph{durghôṭônāy ônek mānush āhôtô hôlô.}
\end{quote}

\begin{quote}
Many people were injured in the accident.
\end{quote}

\begin{cousinweb}
\textbf{\bn{আহত}} (\emph{āhôtô}) is Sanskrit \emph{āhata}, 'struck, beaten, hurt' (Monier-Williams), taken over whole. \bn{আহত} \bn{হওয়া}, to be injured.
\end{cousinweb}

\begin{grammarlens}[title={What you've built}]
One more word for telling what happened.
\end{grammarlens}

\subsection*{Your turn}
\begin{itemize}
\raggedright
  \item \emph{Say it:} \emph{āhôtô}
  \item \emph{Say it:} \emph{āhôtô}, once more
  \item \emph{Say it:} \emph{durghôṭônāy ônek mānush āhôtô hôlô.} — Many people were injured in the accident.
  \item \emph{Recall:} say the Bengali for to break, then the Bengali for to overturn, then say \emph{āhôtô} again
\end{itemize}

\begin{tcolorbox}[breakable,colback=teal!4,colframe=teal!35!black,title={Before you move on}]
What does \bn{আহত} mean? (\textquotedblleft{}Injured; hurt\textquotedblright{}.) Say it once more.
\end{tcolorbox}
